\documentclass[11pt]{article}

\usepackage[margin=1in]{geometry}
\usepackage{amsmath}
\usepackage{enumitem}
\usepackage{array}
\usepackage{parskip}
\usepackage{fancyvrb}

% ---- Fill in your info ----
\newcommand{\myname}{Greyson Knippel}

% Answer block: replace the placeholder text with your answer.
\newenvironment{answer}
  {\par\vspace{0.5em}\noindent\textbf{Answer:}\par\vspace{0.25em}\begingroup\itshape}
  {\endgroup\par\vspace{1em}}

% Blank line for fill-in-the-value questions: \fillin{value}
\newcommand{\fillin}[1]{\underline{\makebox[1.6in][l]{\,#1}}}

\begin{document}

\begin{center}
  {\Large Homework 1 \quad CPSC 346}\\[0.5em]
  \myname
\end{center}

\vspace{1em}

\begin{enumerate}[leftmargin=*]

%===================== Problem 1 =====================
\item The following tables are the current status of the registers and stack
memory for a 64-bit computer. Please explain how the following instructions change the
registers and the stack. Assume a) and b) operate on the original data as shown in the
two tables.

\begin{center}
\begin{tabular}{|l|l|}
\hline
Register & Value \\
\hline
\%rax & 4 \\
\hline
\%rbx & 8 \\
\hline
\%rsp & 0x00FFEE88 \\
\hline
\end{tabular}
\end{center}

\vspace{0.5em}
\noindent\hspace{1em} Memory(Stack):

\begin{center}
\begin{tabular}{|p{1.2in}|p{1.2in}|}
\hline
Address & Data \\
\hline
0x00FFEE88 & 15 \\
\hline
 & \\
\hline
\end{tabular}
\end{center}

\begin{enumerate}[label=\alph*)]
  \item \verb|push %rax|
  \begin{answer}
    if you push \%rax to the stack, \%rsp would point to a location holding 4. and then the stack here would
    grow downward. specifically, \%rsp would go down 8 since its 64 bit, and then \%rsp would be pointing to 0x00FFEE80,
    with the value 4.
  \end{answer}

  \item \verb|pop %rbx|
  \begin{answer}
    after pop \%rbx, becomes 15, and \%rsp is 0x00FFEE90, adding 8 now. 
  \end{answer}
\end{enumerate}

%===================== Problem 2 =====================
\item (10 points) Fill in the value of the registers and the stack memory. Before executing
the instruction ``callq \_Subroutine,'' we have \%rsp = 0xFF08; (\%rsp) = 0x8, which is
the ``top'' value on the stack.

\begin{Verbatim}[xleftmargin=1.5em]
_caller:
   .
   .
   400544: callq  _Subroutine
   400549: mov    %rax,(%rbx)
   .
   .
_Subroutine:
   400650:  mov    %rdi,%rax
   .
   .
   400657:  retq
\end{Verbatim}

\begin{enumerate}[label=\arabic*)]

  \item About to execute the instruction ``callq \_Subroutine'':
  \begin{itemize}[label={},leftmargin=2em]
    \item \%rip: \fillin{}
  \end{itemize}

  \item Right after the instruction ``callq \_Subroutine'' is executed:\\
  (Program will move to the subroutine but not start executing the subroutine yet.)
  \begin{itemize}[label={},leftmargin=2em]
    \item \%rip: \fillin{}
    \item \%rsp: \fillin{}
    \item (\%rsp): \fillin{}
  \end{itemize}

  \item About to execute the instruction ``retq'':
  \begin{itemize}[label={},leftmargin=2em]
    \item \%rip: \fillin{}
    \item \%rsp: \fillin{}
    \item (\%rsp): \fillin{}
  \end{itemize}

  \item Right after the instruction ``retq'' is executed:
  \begin{itemize}[label={},leftmargin=2em]
    \item \%rip: \fillin{}
    \item \%rsp: \fillin{}
    \item (\%rsp): \fillin{}
  \end{itemize}

\end{enumerate}

\end{enumerate}

\end{document}
